\chapter{Network Threat Detection}
\label{ch:3.5}
\chaptertopic{Topic 3.5}

\chapterquote{In God we trust; all others must bring data.}{attributed to W. Edwards Deming}

\section{Motivation}

The firewall from Chapter~\ref{ch:3.4} decides which conversations are allowed to start. This
closing chapter of Unit~3 is about the harder question that follows: does traffic that was
\emph{correctly} permitted still look like an attack? Many real attacks travel over exactly the
ports a firewall must leave open -- HTTP, DNS, encrypted HTTPS. A device that only checks
``is this port allowed'' will never catch a web-application exploit or malware quietly calling
home. This chapter builds the detection layer that sits behind the firewall, then the monitoring
discipline that catches what even detection signatures miss. It closes with the structured
process an organization follows once something is actually found.

\begin{apalignment}{Unit 3: Securing Networks}{Topic 3.5 -- Detecting Network Attacks}
\begin{apobjectives}
\apobjective{3.5.A}{Detect Attacks}{Identify the automated tooling used to detect attacks against a
network -- network and host intrusion detection sensors, inline prevention systems, flow
collectors, full packet capture, SIEM platforms, and network detection and response -- and know
what each can and cannot observe.}
\apobjective{3.5.B}{Detect Attacks}{Explain the structured incident-response lifecycle an
organization follows once an attack is confirmed -- including why containment must precede
eradication, and why capturing lessons learned feeds back into the organization's future
detection and preparation capability.}
\apobjective{3.5.C}{Detect Attacks}{Determine which detection method suits a stated situation --
signature matching against known-bad patterns, statistical or behavioral anomaly detection, or
metadata analysis of encrypted traffic -- and where the sensor must sit to see the traffic at all.}
\apobjective{3.5.D}{Detect Attacks}{Evaluate what a detection method actually costs and delivers:
its precision and recall at real event volumes, the alert load it places on analysts, the traffic
it cannot see, and the operational risk it introduces.}
\apobjective{3.5.E}{Detect Attacks}{Apply detection techniques to log files -- firewall, proxy, DNS,
and flow records -- to recover indicators of an attack such as scan fan-out, periodic beaconing,
tunnelled data, and asymmetric outbound transfer.}
\end{apobjectives}
\end{apalignment}

\section{Intrusion Detection and Prevention Systems}

\begin{definition}{Intrusion Detection System (IDS)}
A passive sensor that inspects a copy of traffic -- via a network tap or port mirror -- against
known attack patterns and raises an alert, without interrupting the traffic itself.
\end{definition}

\begin{definition}{Intrusion Prevention System (IPS)}
An inline device that traffic must physically pass through; on a match it can drop the packet,
reset the connection, or block the source IP in real time, before the payload ever reaches its
target.
\end{definition}

The distinction is deliberately simple: IDS watches and tells, IPS watches and stops. Many
organizations run the identical sensor in both modes on different segments. It runs as IDS on a
high-value segment, where a false positive dropping legitimate traffic would be its own outage,
and as IPS on the perimeter, where automatic blocking is worth the risk.

Two detection philosophies underlie either mode. \textbf{Signature-based detection} matches
traffic against a database of known attack patterns -- a specific SQL-injection string, a malware
file hash, a particular shellcode sequence. It is fast and precise for anything already
catalogued, but blind by construction to a zero-day with no signature yet written.
\textbf{Anomaly-based detection} instead builds a behavioral baseline of what normal traffic looks
like -- typical connection rates, usual protocols per host, expected working-hour activity -- and
flags deviations regardless of whether any signature matches. That lets it catch a genuinely
novel attack, at the cost of a higher false-positive rate: legitimate but unusual behavior
(a 2~a.m.\ backup job, a one-off mass database query) can trigger it just as easily.

\begin{lstlisting}[language=Python, caption={Matching packets against IDS signatures}]
def match_signatures(packets, signatures):
    """packets: list of dicts like {"id": str, "dst_port": int,
       "payload": str}. signatures: list of dicts like
       {"dst_port": int|None, "keyword": str|None}.
       Returns sorted packet IDs matching at least one signature."""
    matched = []
    for p in packets:
        for sig in signatures:
            port_ok = sig["dst_port"] is None or sig["dst_port"] == p["dst_port"]
            keyword_ok = sig["keyword"] is None or sig["keyword"] in p["payload"]
            if port_ok and keyword_ok:
                matched.append(p["id"])
                break
    return sorted(matched)
\end{lstlisting}

Every detection verdict falls into one of four outcomes, and the vocabulary matters because it is
how analysts talk about tuning. A \textbf{true positive} is a real attack correctly caught. A
\textbf{false positive} is benign traffic incorrectly flagged -- the daily burden that produces alert
fatigue. A \textbf{false negative} is a real attack that passed through undetected -- the single most
dangerous outcome. A \textbf{true negative} is benign traffic correctly left alone. A sensor
tuned too sensitively drowns analysts in false positives until they start ignoring alerts
altogether, letting real attacks slip past unnoticed. A sensor tuned too loosely produces false
negatives directly. This tension, not raw detection capability, is what real IDS/IPS management is
actually about.

Sensors also differ in placement. A \textbf{NIDS} (Network IDS) sits at a network choke point and
inspects traffic flowing past. It sees an attack before it reaches any host, but cannot inspect
encrypted payloads without separately performing TLS inspection. A \textbf{HIDS} (Host IDS) runs as
an agent directly on an endpoint, watching system calls, file changes, and process creation. It
sees what actually executed, at the cost of needing an agent on every host it protects. Enterprise
deployments typically run both together.

\begin{misconception}
\textbf{``Signature-based IDS catches all attacks.''} Only attacks with an existing, matching
signature are caught. A zero-day or a sufficiently novel variant passes through undetected unless
anomaly-based detection is also running alongside it.
\end{misconception}

\section{Network Monitoring and Traffic Analysis}

An IDS signature engine catches attacks it already recognizes. Network monitoring is the broader,
complementary discipline of continuously watching traffic itself -- volume, ports, destinations,
timing -- well enough to notice when something looks wrong even with no matching signature at all.

Two views of the same traffic serve different purposes. A full \textbf{packet capture (pcap)}
records every byte of every packet, precise but expensive to store and slow to search at scale.
A \textbf{flow record} (commonly NetFlow) instead summarizes one conversation -- source and
destination IP, ports, protocol, byte and packet counts, duration -- without keeping the payload at
all. That makes it far cheaper to store and query, which is why most enterprise monitoring
pipelines run on flow data, reserving full capture for a smaller set of high-value segments or
triggered deep-dive investigations.

A \textbf{SIEM} (Security Information and Event Management) platform aggregates the alerts and logs
these sensors produce -- IDS/IPS hits, flow records, endpoint events -- into one searchable system.
An analyst is not switching between a dozen separate consoles to piece a timeline together. A
\textbf{network detection and response (NDR)} platform goes one step further: it layers the
anomaly-based monitoring described below directly on top of flow and packet data. It then takes
automated response action itself -- isolating a host, blocking a flow -- collapsing detection and
first-line containment into a single product.

\begin{definition}{Baselining}
Recording typical traffic patterns over time -- which hosts talk to which others, on which ports,
at what volume, during what hours -- to establish what ``normal'' looks like before anything can be
called abnormal.
\end{definition}

Several recurring patterns are worth watching regardless of whether any signature exists for them.
\textbf{Unusual port usage} -- a host suddenly using a port it has never touched before -- is worth
investigating on its own, since the port itself is the anomaly. \textbf{Beaconing} is malware
``phoning home'' to a command-and-control server at suspiciously regular, machine-precise intervals
-- every 60 seconds, every 5 minutes -- rather than the irregular bursts typical of human-driven
traffic. It is exactly the kind of pattern traffic timing analysis catches without ever
inspecting payload content. \textbf{Data-volume spikes} -- a host that normally uploads a few
megabytes suddenly sending several gigabytes externally, especially outside business hours -- are a
strong indicator of exfiltration in progress. The signal here is volume and direction alone, not
what the data actually contains. \textbf{Off-baseline timing} -- an otherwise unremarkable large
database query at 3~a.m.\ from a workstation account -- raises suspicion even when every individual
packet looks entirely benign.

\begin{lstlisting}[language=Python, caption={Flagging hosts that deviate from a per-host traffic baseline}]
def find_anomalous_hosts(baseline, observed):
    """baseline: dict mapping host -> {"volume_threshold": int,
       "allowed_ports": set[int]}.
       observed: list of dicts like {"host": str, "volume": int,
       "port": int}.
       Returns sorted hosts exceeding their volume threshold or using
       a port outside their allowed set."""
    flagged = set()
    for record in observed:
        host = record["host"]
        profile = baseline.get(host)
        if profile is None:
            continue
        if record["volume"] > profile["volume_threshold"]:
            flagged.add(host)
        if record["port"] not in profile["allowed_ports"]:
            flagged.add(host)
    return sorted(flagged)
\end{lstlisting}

A firewall enforces which conversations are allowed to start. An IDS matches traffic against known
bad patterns. Network monitoring catches what neither can: an attack using entirely permitted
ports and protocols, with no matching signature, whose volume, timing, or destination simply
deviates from what the network normally does. This is why mature security operations run all three
layers together rather than trusting any single one.

\begin{misconception}
\textbf{``You need full packet capture to detect an attack.''} Most anomalies -- volume spikes,
beaconing, unusual ports -- are visible in flow records alone, without ever inspecting payload
content at all.
\end{misconception}

\section{Incident Response}

Detecting an attack, by any of the means above, is only the beginning. What an organization does in
the minutes and hours afterward determines whether the incident stays a contained inconvenience or
becomes a catastrophic breach.

\begin{definition}{Incident-handling lifecycle}
The six-phase model popularized by the SANS Institute -- often abbreviated
PICERL\srccite{3.5}{1}: \textbf{Preparation}, \textbf{Identification},
\textbf{Containment}, \textbf{Eradication}, \textbf{Recovery}, and \textbf{Lessons Learned},
followed strictly in that order so that critical steps are never skipped under the pressure of an
active incident. NIST's own guidance, SP~800-61, groups the identical work into four broader
phases instead -- Preparation; Detection and Analysis; Containment, Eradication, and Recovery; and
Post-Incident Activity\srccite{3.5}{2} -- so the two models describe the same underlying process at
different granularity, not two competing processes.
\end{definition}

\begin{figure}[h]
\centering
\begin{tikzpicture}[
  ph/.style={draw=kickstartermaroon, thick, rounded corners, minimum width=2.3cm, minimum height=0.8cm, align=center, font=\scriptsize},
]
\node[ph, fill=kickstartercream] (p1) at (90:2.6) {Preparation};
\node[ph] (p2) at (30:2.6) {Identification};
\node[ph, fill=kickstartermaroon!15] (p3) at (-30:2.6) {Containment};
\node[ph, fill=kickstartermaroon!25] (p4) at (-90:2.6) {Eradication};
\node[ph] (p5) at (-150:2.6) {Recovery};
\node[ph, fill=kickstartercream] (p6) at (150:2.6) {Lessons\\Learned};
\draw[-{Latex[length=2mm]}, kickstartermaroon] (p1) to[bend right=15] (p2);
\draw[-{Latex[length=2mm]}, kickstartermaroon] (p2) to[bend right=15] (p3);
\draw[-{Latex[length=2mm]}, kickstartermaroon] (p3) to[bend right=15] (p4);
\draw[-{Latex[length=2mm]}, kickstartermaroon] (p4) to[bend right=15] (p5);
\draw[-{Latex[length=2mm]}, kickstartermaroon] (p5) to[bend right=15] (p6);
\draw[-{Latex[length=2mm]}, kickstartermaroon, dashed] (p6) to[bend right=15] (p1);
\end{tikzpicture}
\caption{The six-phase (PICERL) incident-handling lifecycle. The dashed return arrow is deliberate:
Lessons Learned feeds directly back into Preparation for the next incident, not a separate,
disconnected step.}
\label{fig:3.5-ir-lifecycle}
\end{figure}

\textbf{Preparation} happens entirely before any incident occurs -- writing the response plan,
defining roles and an escalation chain, deploying the logging and monitoring tools from earlier in
this chapter. It also means running tabletop exercises, so the team is not improvising for the
first time during a real breach. \textbf{Identification} is where an IDS alert or a
traffic-monitoring anomaly actually earns its keep: confirming that an incident is genuinely
occurring rather than routine noise, then scoping which systems, data, and timeframe are affected.
False positives get filtered out here, before resources commit to a full response.

\textbf{Containment} stops the incident from spreading further, without necessarily fixing the root
cause yet -- isolating an infected host, disabling a compromised account, blocking a malicious IP
at the firewall. Practitioners distinguish short-term containment (an immediate, possibly
temporary measure, like unplugging a cable) from long-term containment (a more durable fix that
can hold while eradication is planned). \textbf{Eradication} then removes the actual cause -- deleting
malware, disabling a backdoor account the attacker created, patching the exploited vulnerability.

\begin{example}[Why order is not optional]
Eradicating before containing is a common and costly mistake. Cleaning malware off a host that
remains connected to a live network can tip the attacker off that they have been discovered, or
simply allow immediate reinfection from another already-compromised host on the same segment.
Containment exists specifically to remove that risk before eradication ever begins. The sequence
is not a formality -- it is what prevents cleanup from being undone in real time by an attacker who
is still watching.
\end{example}

\textbf{Recovery} restores systems to normal operation -- rebuilding from clean backups,
reconnecting to the network, closely monitoring the restored systems afterward. It only begins
once eradication is verified complete: bringing a system back online before its root cause is
actually removed just restarts the incident. \textbf{Lessons Learned} closes the loop: documenting a
timeline of what happened, evaluating how well the response plan actually worked, and updating that
plan, tooling, or training based on what was learned. Skipping this phase is how organizations end
up repeating the same incident. The technical fix in Eradication addressed this one instance, but
the process gap that allowed it in the first place -- a missing patch policy, an under-monitored
segment -- goes unaddressed and waits for the next attacker to find it again.

\begin{lstlisting}[language=Python, caption={Validating that a list of incident-response actions follows correct phase order}]
PHASE_ORDER = ["Preparation", "Identification", "Containment",
               "Eradication", "Recovery", "Lessons Learned"]

def validate_sequence(actions):
    """actions: list of dicts like {"description": str, "phase": str}.
       Returns True if phases appear in non-decreasing lifecycle order."""
    last_index = -1
    for action in actions:
        phase_index = PHASE_ORDER.index(action["phase"])
        if phase_index < last_index:
            return False
        last_index = phase_index
    return True
\end{lstlisting}

\begin{misconception}
\textbf{``You should remove the malware as soon as you find it.''} Removing malware before
containment can tip off the attacker or allow immediate reinfection. Contain first, so the
environment is stable, and only then eradicate.
\end{misconception}

\section{Detection Engineering in Practice}

The vocabulary from the previous two sections becomes concrete only once it is written down as a
rule a sensor actually evaluates. This section walks through three rule formats an analyst
encounters directly. It then covers two small calculations -- entropy and timing variance -- that
turn a raw log file into the indicators described above without ever inspecting a payload.

\begin{misconception}
\textbf{``NIST SP 800-94 was updated for the encrypted, cloud-era internet.''} The 2007 edition
(Scarfone \& Mell) remains NIST's only finalized IDS/IPS guidance.\srccite{3.5}{3} A revision was
drafted in 2012 but was later retired \emph{without being finalized} -- NIST itself noted its
inputs no longer matched current threat models. The 2007 document's vocabulary (signature-based,
anomaly-based, stateful-protocol-analysis detection; NIDS, HIDS, wireless IDS) is still the field's
canonical terminology. But the guidance predates encrypted-by-default traffic, cloud workloads, and
machine-learning-based detection entirely, and must be read with that gap in mind -- it was never
brought current, only left in place.
\end{misconception}

A \textbf{Suricata} rule is the canonical format for a signature keyed to packet-level behavior.
The following rule flags a single source IP making an unusually high number of connection attempts
to distinct destination ports in a short window -- a half-open SYN sweep, the classic signature of
active reconnaissance (see Chapter~\ref{ch:3.1}):

\begin{lstlisting}[caption={A Suricata rule detecting a SYN sweep from a single source}]
alert tcp $EXTERNAL_NET any -> $HOME_NET any ( \
    msg:"LOCAL SCAN half-open SYN sweep from single source"; \
    flow:stateless; flags:S,12; \
    threshold:type both, track by_src, count 20, seconds 60; \
    classtype:attempted-recon; \
    reference:url,attack.mitre.org/techniques/T1046/; \
    sid:1000501; rev:1;)
\end{lstlisting}

Each clause does distinct work. The header (\texttt{alert tcp \$EXTERNAL\_NET any ->
\$HOME\_NET any}) states the action, protocol, and direction -- any external host, any port, to any
host and port on the protected network. \texttt{flags:S,12} matches a bare SYN packet with reserved
bits masked, the fingerprint of a half-open scan rather than a normal handshake. The
\texttt{threshold} keyword is what makes this a \emph{sweep} detector rather than a single-packet
rule: it fires only once at least 20 matching packets from the same source arrive within 60
seconds. That converts a noisy raw signal into one alert per incident instead of twenty. The
\texttt{sid} (rule ID) and \texttt{rev} (revision) fields are bookkeeping required by the engine.
\texttt{reference} ties the alert directly to a MITRE ATT\&CK technique, so an analyst can look
up the adjacent techniques an attacker who reached this stage might attempt next.

Not every detector matches packets directly. \textbf{Zeek} (formerly Bro) instead decodes protocols
into structured transaction logs -- one line per connection, DNS query, or TLS handshake -- and
leaves detection to analysis over those records. That lets an analyst ask a question like ``which
internal hosts negotiated TLS with a rare client fingerprint'' instead of ``does any packet match
this byte pattern.'' A short Zeek script expresses exactly that:

\begin{lstlisting}[caption={A Zeek script flagging rare TLS client fingerprints}]
@load base/protocols/ssl
module RareTLSClient;
export {
    redef enum Notice::Type += { Rare_Fingerprint };
    const host_threshold = 2 &redef;
}
global fp_hosts: table[string] of set[addr];

event ssl_established(c: connection)
{
    if ( ! c$ssl?$ja3 ) return;
    local fp = c$ssl$ja3;
    local src = c$id$orig_h;
    if ( fp !in fp_hosts ) fp_hosts[fp] = set();
    add fp_hosts[fp][src];
    if ( |fp_hosts[fp]| <= host_threshold )
        NOTICE([$note=Rare_Fingerprint, $conn=c,
                $msg=fmt("Rare TLS fingerprint %s from %s", fp, src),
                $identifier=cat(fp, src)]);
}
\end{lstlisting}

Organizations change SIEM platforms far more often than they want to rewrite an entire detection
library. \textbf{Sigma} solves that problem: it expresses a log-based detection once, in
vendor-neutral YAML, and a converter compiles it into Splunk SPL, Elasticsearch queries, or another
platform's native syntax. The rule below encodes the DNS-tunnelling indicators from the previous
section directly:

\begin{lstlisting}[caption={A Sigma rule for probable DNS tunnelling via long, high-volume subdomain labels}]
title: Probable DNS Tunnelling by Long High-Volume Subdomain Queries
logsource:
    category: dns_query
detection:
    long_label:
        query|re: '^[A-Za-z0-9+/=_-]{40,63}\.'
    excluded_domains:
        query|endswith:
            - '.in-addr.arpa'
            - '.cloudfront.net'
    timeframe: 10m
    condition: long_label and not excluded_domains | count() by src_ip > 50
\end{lstlisting}

Two of the indicators introduced earlier -- a high-entropy DNS label and a beaconing interval --
are recoverable with nothing more than the Python standard library. \textbf{Shannon entropy}
measures the unpredictability of a string in bits per character. A label like \texttt{api} is
low-entropy (predictable, dictionary-like), while a label carrying base32- or base64-encoded
exfiltrated data is close to the maximum entropy its length allows. The \textbf{coefficient of
variation (CV)} of a series of connection timestamps -- standard deviation divided by the mean gap
-- measures how machine-regular the intervals are. Malware beaconing every 60 seconds produces a
tiny CV, while irregular human browsing produces a large one. Because CV is unitless, a 30-second
beacon and a 30-minute beacon land on the same scale.

\begin{lstlisting}[language=Python, caption={Computing DNS-label entropy and beacon-interval regularity}]
import math
import statistics
from collections import Counter

def shannon_entropy(s):
    """Entropy of s in bits per character (Shannon, 1948)."""
    if not s:
        return 0.0
    n = len(s)
    return -sum((c / n) * math.log2(c / n) for c in Counter(s).values())

def label_score(label):
    """Return (entropy, ceiling) for one DNS label. Entropy cannot
    exceed log2(len(label)), so short labels are capped by arithmetic,
    not content -- absolute entropy paired with length, not a
    normalized ratio, is what separates 'api' from an encoded label."""
    h = shannon_entropy(label)
    ceiling = math.log2(len(label)) if label else 0.0
    return h, ceiling

def interval_cv(timestamps):
    """Coefficient of variation of gaps between sorted timestamps.
    Needs at least three timestamps to give two gaps."""
    if len(timestamps) < 3:
        raise ValueError("need at least three timestamps")
    gaps = [b - a for a, b in zip(timestamps, timestamps[1:])]
    mean = statistics.mean(gaps)
    sd = statistics.pstdev(gaps)
    return mean, sd, (sd / mean if mean else float("inf"))
\end{lstlisting}

Run against a real log, \texttt{shannon\_entropy("api")} returns about 1.58 bits/char (its own
ceiling -- three distinct letters, each equally likely). A longer, DGA-style label such as
\texttt{shannon\_entropy("4a6f68e9")} instead returns close to 4 bits/char, sustained across eight
characters instead of three: short strings can hit their ceiling by chance, but sustained high
entropy over a longer label cannot. Feeding
\texttt{interval\_cv} a fixed 30-minute beacon with 6\% jitter still returns a CV under 0.1, while
ordinary human browsing timestamps return a CV well above 1. The gap is wide enough that a simple
fixed threshold, not a machine-learning model, already separates the two cases correctly.

\begin{misconception}
\textbf{``Normalizing entropy by a label's maximum possible entropy gives a cleaner signal.''} It
does the opposite. A three-letter label like \texttt{api} already sits at its own ceiling
(ratio $1.00$) purely because it is short, while a 26-character encoded label can score a lower
ratio (around $0.79$) despite carrying far more actual disorder. Absolute entropy read alongside
label length is the correct signal -- the normalized ratio actively misleads for short strings.
\end{misconception}

\section*{Review Questions}
\begin{reviewquestions}
\item A newly discovered exploit has no existing signature in an IDS's database. Explain why
signature-based detection misses it entirely and describe the specific property of anomaly-based
detection that could catch it instead, along with the cost that comes with that capability.
\item An IDS is tuned so sensitively that analysts receive hundreds of alerts per day, most of them
benign. Explain, using the TP/FP/FN/TN vocabulary, the specific danger this creates even though the
sensor is technically working as configured.
\item A host that has never used port 4444 begins making regular outbound connections to the same
external IP every 60 seconds, each transferring only a few kilobytes. Explain which two monitoring
indicators from this chapter this scenario exhibits and why payload inspection is not required to
catch it.
\item An incident responder eradicates malware from a compromised host before isolating it from the
network. Explain the specific risk this creates and why containment is sequenced before eradication
in the incident-handling lifecycle.
\item Explain why ``Lessons Learned'' is treated as a required phase of incident response rather
than an optional follow-up, using a concrete example of what goes wrong when it is skipped.
\item The Sigma rule in this chapter's DNS-tunnelling example combines a long-label condition with
an excluded-domains condition rather than flagging every long DNS label. Explain what specific
false-positive source the excluded-domains clause is guarding against, and why omitting it would
degrade the rule's precision without improving its recall.
\item A 30-minute beacon with 6\% timing jitter still produces a coefficient of variation under
0.1, clearly separable from ordinary human browsing. Explain why CV, rather than the raw standard
deviation of the intervals, is the right statistic to compare a 30-second beacon against a
30-minute beacon on the same scale.
\end{reviewquestions}

\begin{answerkey}
\item Signature-based detection matches traffic only against a database of already-catalogued
patterns, so an exploit with no existing signature simply matches nothing and passes through
undetected by construction, regardless of how obviously malicious its behavior actually is.
Anomaly-based detection instead builds a behavioral baseline of normal activity and flags any
significant deviation from it. That lets it catch a genuinely novel attack with no signature at
all -- at the cost of a higher false-positive rate, since legitimate but unusual behavior can
deviate from the baseline just as sharply as an actual attack does.
\answerlabel{3.5.C}
\item Hundreds of daily alerts, most of them false positives, produce alert fatigue: analysts begin
tuning out or triaging alerts less carefully simply to keep pace with the volume. The specific
danger is that this raises the effective false-negative rate even though the sensor itself is
technically functioning as configured. A real attack (a true positive) can arrive buried among the
noise and be dismissed or missed entirely by an exhausted analyst -- a failure of the human
response process sitting downstream of a technically correct detection.
\answerlabel{3.5.D}
\item This scenario exhibits unusual port usage (a host using port 4444, which it has never touched
before) and beaconing (regular, machine-precise 60-second intervals to the same external IP,
transferring a near-constant small volume each time). That interval pattern is exactly the
coefficient-of-variation signature distinguishing malware beaconing from irregular human-driven
traffic. Payload inspection is not required because both indicators live entirely in metadata --
which port is used, how regular the timing is, how much data moves -- all of which is visible in
flow records even when the payload itself is fully encrypted.
\answerlabel{3.5.E}
\item Eradicating malware from a host that remains connected to the network risks tipping off the
attacker that they have been discovered, prompting them to escalate or destroy evidence. It can
also allow immediate reinfection from another already-compromised host on the same segment that was
never isolated. Containment is sequenced before eradication in this lifecycle specifically to
remove that risk first -- stabilizing the environment so that cleanup performed during eradication
cannot be undone in real time by an attacker who is still actively watching or still has other
footholds available.
\answerlabel{3.5.A}
\item Lessons Learned is required, not optional, because Eradication only fixes the specific
technical instance of an incident -- the malware, the exploited vulnerability -- while leaving
unaddressed whatever process gap allowed the incident to happen in the first place, such as a
missing patch policy or an under-monitored network segment. If this phase is skipped, the
organization documents no timeline, evaluates nothing about how well the response plan performed,
and updates neither tooling nor training. That is precisely how the same category of incident
recurs: the next attacker finds the identical unaddressed gap the first one exploited.
\answerlabel{3.5.B}
\item Content-delivery and reverse-DNS infrastructure legitimately generates long, high-entropy-
looking subdomain labels as a normal part of how those services operate -- a CDN edge node or a
reverse-lookup zone is not evidence of tunnelling. The excluded-domains clause filters out exactly
this known-benign source before the count threshold is evaluated. Omitting it would not catch any
additional real tunnelling traffic that the long-label condition was not already going to catch.
It would only add legitimate CDN and reverse-DNS traffic to the alert population, which lowers
precision (more false positives) while leaving recall (the fraction of real tunnelling caught)
completely unchanged.
\answerlabel{3.5.E}
\item Standard deviation alone is expressed in the same units as the interval itself, so a
30-second beacon's jitter and a 30-minute beacon's jitter produce very different raw
standard-deviation numbers even at the identical 6\% relative regularity. That makes the two
incomparable on one fixed threshold. Dividing by the mean gap (producing CV) normalizes away the
absolute timescale, so a beacon's \emph{relative} regularity -- how small its jitter is compared to
its own interval -- is what gets compared. Both a 30-second and a 30-minute beacon at the same
relative regularity land on the same CV scale, which is exactly the property that lets one fixed
threshold work regardless of the beacon's period.
\answerlabel{3.5.E}
\end{answerkey}

\section*{Sources}
\begin{enumerate}[label=\lbrack\arabic*\rbrack, leftmargin=2.2em]
\item \srcitem{3.5}{1}Patrick Kral, \emph{The Incident Handler's Handbook}, SANS Institute
InfoSec Reading Room, 2011 -- the widely-taught source for the six-phase PICERL (Preparation,
Identification, Containment, Eradication, Recovery, Lessons Learned) incident-handling model.
\item \srcitem{3.5}{2}Paul Cichonski, Tom Millar, Tim Grance, and Karen Scarfone, \emph{Computer
Security Incident Handling Guide}, NIST Special Publication 800-61 Revision 2, National Institute
of Standards and Technology, August 2012.
\item \srcitem{3.5}{3}Karen Scarfone and Peter Mell, \emph{Guide to Intrusion Detection and
Prevention Systems (IDPS)}, NIST Special Publication 800-94, National Institute of Standards and
Technology, February 2007; Draft NIST SP 800-94 Revision 1 (July 2012) was never finalized and was
formally retired in 2022.
\end{enumerate}
